\section{Síntesis y ampliación}

\begin{table}[H]
\centering
\begin{tabular}{p{3.5cm}p{8cm}}
\toprule
Dimensión & Revisión central \\
\midrule
Motivación & La estructura jerárquica es fundamental para resolver la asignación de crédito y la exploración \\
Métodos & Options/SMDP + goal-conditioned + imitation (VQ-VAE) ofrecen una ruta completa \\
Gobernanza del lenguaje & LLM prompt log + verification pipeline garantizan que el plan de alto nivel sea confiable \\
Ingeniería | evaluación & scoreboard + incident grade-book permiten rastrear el emergent drift \\
\bottomrule
\end{tabular}
\caption{Repaso de las tres líneas principales de la Lecture 15.}
\end{table}

\begin{table}[H]
\centering
\begin{tabular}{p{3.5cm}p{8cm}}
\toprule
Dimensión & Riesgos y control \\
\midrule
No estacionariedad & hindsight relabel + off-policy correction \\
Hallucination & multi-metric governance + tool check \\
Script drift & versioned prompt log + incident reviews \\
\bottomrule
\end{tabular}
\caption{Perspectiva complementaria sobre riesgos y gobernanza.}
\end{table}

\subsection{Lecturas adicionales}
\begin{enumerate}
    \item Sutton et al., ``Between MDPs and Semi-MDPs: A Framework for Temporal Abstraction in RL,'' 1999
    \item Nachum et al., ``Data-Efficient Hierarchical Reinforcement Learning (HIRO),'' NeurIPS 2018
    \item Pertsch et al., ``Accelerating RL with Learned Skill Priors (SPiRL),'' CoRL 2020
    \item Ahn et al., ``Do As I Can, Not As I Say: Grounding Language in Robotic Affordances (SayCan),'' 2022
    \item Ajay et al., ``Compositional Foundation Models for Hierarchical Planning,'' NeurIPS 2023
\end{enumerate}

\end{document}
